% Generated draft: requires mathematical and source review.
\documentclass[11pt]{article}
\usepackage[margin=25mm]{geometry}
\usepackage{fontspec}
\setmainfont{Latin Modern Roman}
\usepackage{amsmath,amssymb,mathtools,unicode-math}
\setmathfont{Latin Modern Math}
\usepackage{enumitem,needspace,xurl,hyperref}
\setlength{\parindent}{0pt}
\setlength{\parskip}{5pt}
\setlength{\emergencystretch}{3em}
\newcommand{\N}{\mathbb{N}}
\newcommand{\Z}{\mathbb{Z}}
\newcommand{\Q}{\mathbb{Q}}
\newcommand{\R}{\mathbb{R}}
\newcommand{\C}{\mathbb{C}}
\newcommand{\F}{\mathbb{F}}
\newcommand{\A}{\mathbb{A}}
\newcommand{\PP}{\mathbb{P}}
\newcommand{\E}{\mathbb{E}}
\newcommand{\eps}{\varepsilon}
\newcommand{\dd}{\,\mathrm{d}}
\begin{document}
\section*{UnsolvedMath ID: 30004328}
\textbf{Source problem code:} catalogue entry
\subsection{Definitions and mathematical statement}
Let \(X\) be a smooth projective surface over a field \(K\). The supplied source specifies neither \(K\) nor its characteristic or algebraic closure assumptions. For Cartier divisors \(D,D'\) on \(X\), write \(D\equiv_{\mathrm{num}}D'\) if \(D\cdot C=D'\cdot C\) for every integral curve \(C\subset X\). Set
\[
L_X=\operatorname{Div}_{\mathrm{Cart}}(X)/\!\equiv_{\mathrm{num}}.
\]
For the following numerical-class formulation, a class \(F\in L_X\) is called primitive if \(F\ne0\) and there are no \(G\in L_X\) and integer \(m\ge2\) with \(F=mG\). An effective divisor is a finite sum of integral curves with nonnegative integer coefficients. Define the semi-effective order, when the displayed set is nonempty, by
\[
s_X(F)=\min\bigl\{k\in\Z_{>0}:\text{there exists an effective divisor }D\text{ on }X\text{ with }[D]=kF\text{ in }L_X\bigr\}.
\]
The question is whether there exist a single smooth projective surface \(X\) and a sequence \((F_n)_{n\ge1}\) of primitive classes in \(L_X\), each having finite semi-effective order \(k_n=s_X(F_n)\), such that
\[
\lim_{n\to\infty}k_n=+\infty.
\]
Thus the surface is fixed throughout the sequence, and for every integer \(M\ge1\) there must be an \(N\ge1\) such that \(k_n>M\) for all \(n\ge N\).

The accompanying irreducible-curve question has the following natural strengthened reading: can such an \(X\) and sequence be chosen with integral curves \(C_n\subset X\) satisfying
\[
[C_n]=k_nF_n,\qquad k_n=s_X(F_n),\qquad k_n\longrightarrow+\infty?
\]
Here an integral curve is taken with coefficient one, so it is a prime divisor. The source does not explicitly say whether irreducibility must hold for a divisor realizing the least effective multiple or merely for some effective multiple; the displayed strengthened reading requires confirmation. The source also does not explicitly identify the integral group in which primitivity is intended. The use of \(L_X\) above is an explicit numerical-class interpretation requiring review, and no assertion about geometric irreducibility is intended without further specification of the base field.
\subsection{Short English statement}
Can one fixed smooth projective surface carry primitive integral numerical divisor classes whose least positive effective multiples have orders tending to infinity? Can this also happen when those least multiples are represented by irreducible, reduced curves? The base field, the group governing primitivity, and the precise irreducibility requirement need confirmation from the source.
\subsection{Sources}
\begin{itemize}
\item Supplied local UnsolvedMath registry record, permanent ID \(30004328\), title ``Unbounded Semi Effective Orders on Projective Surfaces,'' source problem code \texttt{catalogue entry}, in the Oberwolfach Reports Open Problems collection. The record gives proposed year \(2019\), names no individual proposer, and labels the problem open. Its \texttt{corrected\_verified} statement label is accompanied by an explanation that only formula delimiters, notation, grammar, and typography were repaired; it does not establish mathematical correctness.
\item Supplied source citation: \emph{Mini-Workshop: Seshadri Constants (2020)}, Oberwolfach Reports, pp.\ 3271--3313, DOI \texttt{10.4171/owr/2019/53}. \url{https://doi.org/10.4171/owr/2019/53}. The supplied background describes Oberwolfach Reports as documenting mathematical workshops and their open-problem discussions.
\item Supplied literature reference: Bauer--Harbourne--Kuronya--Nickel, \emph{Bounded volume denominators and bounded negativity}. \url{https://arxiv.org/abs/1904.07486}.
\item Supplied literature reference: OWR report on Seshadri constants. \url{https://ems.press/content/serial-article-files/46833}.
\item Literature assessment dated exactly \texttt{2026-08-22}: the supplied curation labels the problem open and states that the prime-divisor case with semi-effective order greater than two remains unresolved, whereas higher orders are known for non-prime or reducible effective multiples. This is supplied dated curation derived from the OWR material, not an independently verified literature assessment. The record itself requests independent verification of time-sensitive status and cited claims.
\end{itemize}
\end{document}
